\SourceLeaf{220}
\Needspace{12\baselineskip}
\begin{center}
\textit{General solution in integers of all indeterminate equations of the second degree in two unknowns.}
\end{center}

As a finishing touch to these investigations, we add

\Needspace{8\baselineskip}
\Article{216}
\textsc{Problem.} \textit{Find all integer solutions of the general}\HistoricalNote{*)}{If an equation is given whose second, fourth or fifth coefficient is not even, multiplying it by 2 gives it the form assumed here.} \textit{indeterminate equation of the second degree in two unknowns}
\[
axx+2bxy+cyy+2dx+2ey+f=0
\]
\textit{(where $a,b,c$ etc. are any given integers).}

\textit{Solution.} In place of the unknowns $x,y$, introduce
\[
p=(bb-ac)x+be-cd\quad\text{and}\quad q=(bb-ac)y+bd-ae
\]
which are clearly always integers when $x,y$ are integers. This gives the equation
\[
app+2bpq+cqq+f(bb-ac)^2+(bb-ac)(aee-2bde+cdd)=0
\]
or, putting for brevity
\[
f(bb-ac)^2+(bb-ac)(aee-2bde+cdd)=-M
\]
the equation
\[
app+2bpq+cqq=M
\]
We have already shown how to find all solutions of this equation, \textit{i.e.}, all representations of $M$ by $(a,b,c)$. From each set of values of $p,q$, determine the corresponding values of $x,y$ by the equations
\[
x=\frac{p+cd-be}{bb-ac},\quad y=\frac{q+ae-bd}{bb-ac}
\]
We easily see that all these values satisfy the given equation, and that no integer values of $x,y$ are missed by this method. So if we discard the fractional values from all the values of $x,y$ obtained in this way, all the required solutions remain.

We make the following observations about this solution.

\SourceLeaf{221}
$1^{\mathrm{o}}$ If either $M$ cannot be represented by $(a,b,c)$, or no representation gives \textit{integer} values of $x,y$, then the equation has no integer solution.

$2^{\mathrm{o}}$ If the determinant of $(a,b,c)$, \textit{i.e.}, $bb-ac$, is negative, or is a positive square and at the same time $M$ is not $=0$, the number of representations of $M$ by $(a,b,c)$ is finite. So the number of all solutions of the given equation, if there are any, is also finite.

$3^{\mathrm{o}}$ If $bb-ac$ is positive and nonsquare, or is a square and at the same time $M=0$, then $M$, if it can be represented at all, can be represented in \textit{infinitely many different ways} by $(a,b,c)$. But it is impossible to find all these representations \textit{themselves} and try whether they give integer or fractional values of $x,y$. So we need a rule that lets us be \textit{certain} when no representation at all can give integer values of $x,y$ (for however many representations we \textit{try} in this case, without such a rule we would never reach certainty). If some representations give integer values of $x,y$ and others fractional values, we must show how to distinguish the latter from the former in advance, in general.

$4^{\mathrm{o}}$ If $bb-ac=0$, the values of $x,y$ cannot be determined at all by the preceding formulas. So we must find a \textit{special method} for this case.

